% residual_block.tex
% Detailná štruktúra jedného reziduálneho bloku modelu EdgeResMLP
% Použitie: \resizebox{0.55\textwidth}{!}{% residual_block.tex
% Detailná štruktúra jedného reziduálneho bloku modelu EdgeResMLP
% Použitie: \resizebox{0.55\textwidth}{!}{% residual_block.tex
% Detailná štruktúra jedného reziduálneho bloku modelu EdgeResMLP
% Použitie: \resizebox{0.55\textwidth}{!}{% residual_block.tex
% Detailná štruktúra jedného reziduálneho bloku modelu EdgeResMLP
% Použitie: \resizebox{0.55\textwidth}{!}{\input{figures/tikz/residual_block}}
\begin{tikzpicture}[
    font=\small,
    layer/.style={
        draw, rounded corners=3pt, align=center,
        minimum width=3.2cm, minimum height=0.68cm,
        inner sep=3pt, fill=green!8, draw=green!50
    },
    io/.style={
        draw, rounded corners=3pt, align=center,
        minimum width=3.2cm, minimum height=0.65cm,
        inner sep=3pt, fill=gray!12, draw=gray!50
    },
    sumnode/.style={
        circle, draw=blue!50, fill=blue!6,
        minimum size=0.55cm, inner sep=1pt,
        font=\normalsize\bfseries
    },
    arrow/.style={-{Latex[length=2mm]}, semithick},
    resarrow/.style={-{Latex[length=2mm]}, semithick, draw=green!60!black, thick},
    node distance=0.52cm
]

%% ── Vstup ────────────────────────────────────────────────────────
\node[io]    (inp)    at (1.8, 0)    {\(h_{ij}^{(\ell)}\) (vstup bloku)};

%% ── Priama cesta (hlavná vetva) ──────────────────────────────────
\node[layer] (ln)     at (1.8, -1.10) {LayerNorm};
\node[layer] (lin1)   at (1.8, -2.10) {Linear \(\to d_h\)};
\node[layer] (gelu)   at (1.8, -3.10) {GELU};
\node[layer] (drop)   at (1.8, -4.10) {Dropout(\(p\))};
\node[layer] (lin2)   at (1.8, -5.10) {Linear \(\to d_h\)};

%% ── Sčítanie (reziduálna skratka + výstup transformácie) ─────────
\node[sumnode] (plus) at (1.8, -6.30) {\(+\)};

%% ── Výstup ───────────────────────────────────────────────────────
\node[io]    (out)    at (1.8, -7.30) {\(h_{ij}^{(\ell+1)}\) (výstup bloku)};

%% ── Šípky hlavnej vetvy ──────────────────────────────────────────
\draw[arrow] (inp)  -- (ln);
\draw[arrow] (ln)   -- (lin1);
\draw[arrow] (lin1) -- (gelu);
\draw[arrow] (gelu) -- (drop);
\draw[arrow] (drop) -- (lin2);
\draw[arrow] (lin2) -- (plus);
\draw[arrow] (plus) -- (out);

%% ── Reziduálna skratka (vpravo) ──────────────────────────────────
%% Vetva z inp.east → zakrivenie dole → plus.east
\draw[resarrow]
    (inp.east) -- ++(1.1, 0)
               -- ++(0, -6.30)
               -- (plus.east);

%% Popis skratky
\node[font=\scriptsize, green!60!black, rotate=90, anchor=south]
    at (3.35, -3.15) {reziduálna skratka};

%% ── Popis blokovej štruktúry ─────────────────────────────────────
\node[font=\scriptsize, gray, align=center] at (1.8, -8.10)
    {\(F_\ell(x) = \operatorname{Dropout}(W_2 \cdot \operatorname{GELU}(\operatorname{Dropout}(W_1 x)))\)};

%% Orámovanie celého bloku
\draw[draw=green!40, rounded corners=8pt, dashed]
    (-0.05, 0.40) rectangle (3.65, -5.55);
\node[font=\scriptsize\itshape, green!60!black] at (1.8, 0.55)
    {Reziduálny blok \(\ell\)};

\end{tikzpicture}
}
\begin{tikzpicture}[
    font=\small,
    layer/.style={
        draw, rounded corners=3pt, align=center,
        minimum width=3.2cm, minimum height=0.68cm,
        inner sep=3pt, fill=green!8, draw=green!50
    },
    io/.style={
        draw, rounded corners=3pt, align=center,
        minimum width=3.2cm, minimum height=0.65cm,
        inner sep=3pt, fill=gray!12, draw=gray!50
    },
    sumnode/.style={
        circle, draw=blue!50, fill=blue!6,
        minimum size=0.55cm, inner sep=1pt,
        font=\normalsize\bfseries
    },
    arrow/.style={-{Latex[length=2mm]}, semithick},
    resarrow/.style={-{Latex[length=2mm]}, semithick, draw=green!60!black, thick},
    node distance=0.52cm
]

%% ── Vstup ────────────────────────────────────────────────────────
\node[io]    (inp)    at (1.8, 0)    {\(h_{ij}^{(\ell)}\) (vstup bloku)};

%% ── Priama cesta (hlavná vetva) ──────────────────────────────────
\node[layer] (ln)     at (1.8, -1.10) {LayerNorm};
\node[layer] (lin1)   at (1.8, -2.10) {Linear \(\to d_h\)};
\node[layer] (gelu)   at (1.8, -3.10) {GELU};
\node[layer] (drop)   at (1.8, -4.10) {Dropout(\(p\))};
\node[layer] (lin2)   at (1.8, -5.10) {Linear \(\to d_h\)};

%% ── Sčítanie (reziduálna skratka + výstup transformácie) ─────────
\node[sumnode] (plus) at (1.8, -6.30) {\(+\)};

%% ── Výstup ───────────────────────────────────────────────────────
\node[io]    (out)    at (1.8, -7.30) {\(h_{ij}^{(\ell+1)}\) (výstup bloku)};

%% ── Šípky hlavnej vetvy ──────────────────────────────────────────
\draw[arrow] (inp)  -- (ln);
\draw[arrow] (ln)   -- (lin1);
\draw[arrow] (lin1) -- (gelu);
\draw[arrow] (gelu) -- (drop);
\draw[arrow] (drop) -- (lin2);
\draw[arrow] (lin2) -- (plus);
\draw[arrow] (plus) -- (out);

%% ── Reziduálna skratka (vpravo) ──────────────────────────────────
%% Vetva z inp.east → zakrivenie dole → plus.east
\draw[resarrow]
    (inp.east) -- ++(1.1, 0)
               -- ++(0, -6.30)
               -- (plus.east);

%% Popis skratky
\node[font=\scriptsize, green!60!black, rotate=90, anchor=south]
    at (3.35, -3.15) {reziduálna skratka};

%% ── Popis blokovej štruktúry ─────────────────────────────────────
\node[font=\scriptsize, gray, align=center] at (1.8, -8.10)
    {\(F_\ell(x) = \operatorname{Dropout}(W_2 \cdot \operatorname{GELU}(\operatorname{Dropout}(W_1 x)))\)};

%% Orámovanie celého bloku
\draw[draw=green!40, rounded corners=8pt, dashed]
    (-0.05, 0.40) rectangle (3.65, -5.55);
\node[font=\scriptsize\itshape, green!60!black] at (1.8, 0.55)
    {Reziduálny blok \(\ell\)};

\end{tikzpicture}
}
\begin{tikzpicture}[
    font=\small,
    layer/.style={
        draw, rounded corners=3pt, align=center,
        minimum width=3.2cm, minimum height=0.68cm,
        inner sep=3pt, fill=green!8, draw=green!50
    },
    io/.style={
        draw, rounded corners=3pt, align=center,
        minimum width=3.2cm, minimum height=0.65cm,
        inner sep=3pt, fill=gray!12, draw=gray!50
    },
    sumnode/.style={
        circle, draw=blue!50, fill=blue!6,
        minimum size=0.55cm, inner sep=1pt,
        font=\normalsize\bfseries
    },
    arrow/.style={-{Latex[length=2mm]}, semithick},
    resarrow/.style={-{Latex[length=2mm]}, semithick, draw=green!60!black, thick},
    node distance=0.52cm
]

%% ── Vstup ────────────────────────────────────────────────────────
\node[io]    (inp)    at (1.8, 0)    {\(h_{ij}^{(\ell)}\) (vstup bloku)};

%% ── Priama cesta (hlavná vetva) ──────────────────────────────────
\node[layer] (ln)     at (1.8, -1.10) {LayerNorm};
\node[layer] (lin1)   at (1.8, -2.10) {Linear \(\to d_h\)};
\node[layer] (gelu)   at (1.8, -3.10) {GELU};
\node[layer] (drop)   at (1.8, -4.10) {Dropout(\(p\))};
\node[layer] (lin2)   at (1.8, -5.10) {Linear \(\to d_h\)};

%% ── Sčítanie (reziduálna skratka + výstup transformácie) ─────────
\node[sumnode] (plus) at (1.8, -6.30) {\(+\)};

%% ── Výstup ───────────────────────────────────────────────────────
\node[io]    (out)    at (1.8, -7.30) {\(h_{ij}^{(\ell+1)}\) (výstup bloku)};

%% ── Šípky hlavnej vetvy ──────────────────────────────────────────
\draw[arrow] (inp)  -- (ln);
\draw[arrow] (ln)   -- (lin1);
\draw[arrow] (lin1) -- (gelu);
\draw[arrow] (gelu) -- (drop);
\draw[arrow] (drop) -- (lin2);
\draw[arrow] (lin2) -- (plus);
\draw[arrow] (plus) -- (out);

%% ── Reziduálna skratka (vpravo) ──────────────────────────────────
%% Vetva z inp.east → zakrivenie dole → plus.east
\draw[resarrow]
    (inp.east) -- ++(1.1, 0)
               -- ++(0, -6.30)
               -- (plus.east);

%% Popis skratky
\node[font=\scriptsize, green!60!black, rotate=90, anchor=south]
    at (3.35, -3.15) {reziduálna skratka};

%% ── Popis blokovej štruktúry ─────────────────────────────────────
\node[font=\scriptsize, gray, align=center] at (1.8, -8.10)
    {\(F_\ell(x) = \operatorname{Dropout}(W_2 \cdot \operatorname{GELU}(\operatorname{Dropout}(W_1 x)))\)};

%% Orámovanie celého bloku
\draw[draw=green!40, rounded corners=8pt, dashed]
    (-0.05, 0.40) rectangle (3.65, -5.55);
\node[font=\scriptsize\itshape, green!60!black] at (1.8, 0.55)
    {Reziduálny blok \(\ell\)};

\end{tikzpicture}
}
\begin{tikzpicture}[
    font=\small,
    layer/.style={
        draw, rounded corners=3pt, align=center,
        minimum width=3.2cm, minimum height=0.68cm,
        inner sep=3pt, fill=green!8, draw=green!50
    },
    io/.style={
        draw, rounded corners=3pt, align=center,
        minimum width=3.2cm, minimum height=0.65cm,
        inner sep=3pt, fill=gray!12, draw=gray!50
    },
    sumnode/.style={
        circle, draw=blue!50, fill=blue!6,
        minimum size=0.55cm, inner sep=1pt,
        font=\normalsize\bfseries
    },
    arrow/.style={-{Latex[length=2mm]}, semithick},
    resarrow/.style={-{Latex[length=2mm]}, semithick, draw=green!60!black, thick},
    node distance=0.52cm
]

%% ── Vstup ────────────────────────────────────────────────────────
\node[io]    (inp)    at (1.8, 0)    {\(h_{ij}^{(\ell)}\) (vstup bloku)};

%% ── Priama cesta (hlavná vetva) ──────────────────────────────────
\node[layer] (ln)     at (1.8, -1.10) {LayerNorm};
\node[layer] (lin1)   at (1.8, -2.10) {Linear \(\to d_h\)};
\node[layer] (gelu)   at (1.8, -3.10) {GELU};
\node[layer] (drop)   at (1.8, -4.10) {Dropout(\(p\))};
\node[layer] (lin2)   at (1.8, -5.10) {Linear \(\to d_h\)};

%% ── Sčítanie (reziduálna skratka + výstup transformácie) ─────────
\node[sumnode] (plus) at (1.8, -6.30) {\(+\)};

%% ── Výstup ───────────────────────────────────────────────────────
\node[io]    (out)    at (1.8, -7.30) {\(h_{ij}^{(\ell+1)}\) (výstup bloku)};

%% ── Šípky hlavnej vetvy ──────────────────────────────────────────
\draw[arrow] (inp)  -- (ln);
\draw[arrow] (ln)   -- (lin1);
\draw[arrow] (lin1) -- (gelu);
\draw[arrow] (gelu) -- (drop);
\draw[arrow] (drop) -- (lin2);
\draw[arrow] (lin2) -- (plus);
\draw[arrow] (plus) -- (out);

%% ── Reziduálna skratka (vpravo) ──────────────────────────────────
%% Vetva z inp.east → zakrivenie dole → plus.east
\draw[resarrow]
    (inp.east) -- ++(1.1, 0)
               -- ++(0, -6.30)
               -- (plus.east);

%% Popis skratky
\node[font=\scriptsize, green!60!black, rotate=90, anchor=south]
    at (3.35, -3.15) {reziduálna skratka};

%% ── Popis blokovej štruktúry ─────────────────────────────────────
\node[font=\scriptsize, gray, align=center] at (1.8, -8.10)
    {\(F_\ell(x) = \operatorname{Dropout}(W_2 \cdot \operatorname{GELU}(\operatorname{Dropout}(W_1 x)))\)};

%% Orámovanie celého bloku
\draw[draw=green!40, rounded corners=8pt, dashed]
    (-0.05, 0.40) rectangle (3.65, -5.55);
\node[font=\scriptsize\itshape, green!60!black] at (1.8, 0.55)
    {Reziduálny blok \(\ell\)};

\end{tikzpicture}
